\documentclass[a4paper]{article}

% Español (México) con XeLaTeX: fontspec para fuentes Unicode, polyglossia
% para la división silábica y los nombres del documento en la variante mexicana.
\usepackage{fontspec}
\usepackage{polyglossia}
\setdefaultlanguage[variant=mexican]{spanish}
% Los nombres chinos van como «pinyin (original)»: xeCJK da a los caracteres
% Han su propia fuente; sin ella XeLaTeX los omite con un aviso.
% FandolSong no cubre algunos caracteres de nombres (U+73FA); WenQuanYi Zen Hei sí.
\usepackage[AutoFallBack=true]{xeCJK}
\setCJKmainfont{FandolSong-Regular.otf}
\setCJKfallbackfamilyfont{\CJKrmdefault}{WenQuanYi Zen Hei}
% polyglossia no traduce los nombres de listings.
\AtBeginDocument{\providecommand{\lstlistingname}{}\renewcommand{\lstlistingname}{Listado}%
  \providecommand{\lstlistlistingname}{}\renewcommand{\lstlistlistingname}{Índice de listados}}
\usepackage{amsmath, amssymb}
\usepackage{graphicx}
\usepackage[margin=2.5cm]{geometry}
\usepackage[most]{tcolorbox}
\usepackage{etoolbox}
\usepackage{listings}
\usepackage{hyperref}
\usepackage{booktabs}
\usepackage{subcaption}
\usepackage{float}

\newtcolorbox{knowledgebox}[1]{
    enhanced, colback=blue!5!white, colframe=blue!75!black, colbacktitle=blue!75!black,
    coltitle=white, fonttitle=\bfseries, title={#1}, attach boxed title to top left={yshift=-2mm, xshift=2mm},
    boxrule=1pt, sharp corners
}
\newtcolorbox{importantbox}[1]{
    enhanced, colback=yellow!10!white, colframe=yellow!80!black, colbacktitle=yellow!80!black,
    coltitle=black, fonttitle=\bfseries, title={#1}, sharp corners
}
\newtcolorbox{warningbox}[1]{
    enhanced, colback=red!5!white, colframe=red!75!black, colbacktitle=red!75!black,
    coltitle=white, fonttitle=\bfseries, title={#1}, sharp corners
}

\lstset{
    language=Python, basicstyle=\ttfamily\small, keywordstyle=\color{blue},
    stringstyle=\color{red!60!black}, commentstyle=\color{green!60!black},
    breaklines=true, frame=single, numbers=left, numberstyle=\tiny\color{gray},
    captionpos=b, extendedchars=true
}

\newcommand{\notetitle}{[LLM Agents SP25] Abstraction \& Discovery with LLM Agents — Swarat Chaudhuri}
\newcommand{\noteauthors}{Elaboradas a partir de materiales públicos del curso}
\newcommand{\notedate}{2025}
\newcommand{\videochannel}{Berkeley RDI}
\newcommand{\videopublishdate}{2025}
\newcommand{\videoduration}{aproximadamente 60 minutos}
\newcommand{\videourl}{}
\newcommand{\videocoverpath}{cover.jpg}
\newcommand{\repourl}{https://github.com/hqhq1025/ai-course-notes}


% Footer with repo info
\usepackage{fancyhdr}
\pagestyle{fancy}
\fancyhf{}
\fancyfoot[C]{\thepage}
\fancyfoot[R]{\footnotesize\color{gray}\href{\repourl}{AI Course Notes}}
\renewcommand{\headrulewidth}{0pt}
\renewcommand{\footrulewidth}{0pt}

\begin{document}

\begin{titlepage}
\centering
{\Large Notas del curso\par}
\vspace{1.2cm}
{\huge\bfseries \notetitle\par}
\vspace{0.8cm}
{\large \noteauthors\par}
\vspace{0.3cm}
{\large \notedate\par}
\vspace{1.2cm}
\ifdefempty{\videocoverpath}{}{\includegraphics[width=0.82\textwidth,height=0.45\textheight,keepaspectratio]{\videocoverpath}\par}
\vfill
\begin{tcolorbox}[width=0.9\textwidth, colback=black!2!white, colframe=black!60, sharp corners]
\textbf{Autor o canal del video}: \videochannel\par
\textbf{Fecha de publicación}: \videopublishdate\par
\textbf{Duración del video}: \videoduration\par
\ifdefempty{\videourl}{}{\textbf{Enlace del video}: \href{\videourl}{\nolinkurl{\videourl}}\par}\par
\textbf{Página del proyecto}: \href{\repourl}{\nolinkurl{\repourl}}\par
\end{tcolorbox}
\end{titlepage}

\tableofcontents
\newpage

\section{Introducción: LLM como herramienta para el descubrimiento científico y matemático}

Esta clase está impartida por el profesor de UT Austin Swarat Chaudhuri, y se centra en presentar las aplicaciones de los LLM en el descubrimiento matemático y el descubrimiento científico, con énfasis particular en cómo la capacidad de abstracción de los LLM acelera el proceso de descubrimiento.

\begin{knowledgebox}{Proceso del descubrimiento matemático y del descubrimiento científico}
\textbf{Descubrimiento matemático}: modelado $\to$ conjetura $\to$ razonamiento riguroso (demostración) $\to$ contraejemplo/nueva conjetura $\to$ ciclo iterativo.\\
\textbf{Descubrimiento científico}: construcción de teoría $\to$ planteamiento de hipótesis $\to$ diseño experimental $\to$ verificación experimental $\to$ nuevos datos $\to$ nueva teoría.\\
AI for Math busca automatizar el proceso de demostración; AI for Science busca automatizar la generación de hipótesis y el diseño experimental.
\end{knowledgebox}

\subsection{Las tres ventajas principales del LLM Agent en el descubrimiento}

\begin{enumerate}
    \item \textbf{Orientación de la búsqueda}: los LLM contienen internamente una gran cantidad de conocimiento previo, que puede orientar la dirección de la búsqueda en un espacio de hipótesis enorme
    \item \textbf{Aprendizaje por experiencia}: el LLM Agent puede interactuar con el entorno y aprender de la experiencia (incorporándola al prompt o ajustando los pesos mediante RL)
    \item \textbf{Capacidad de abstracción}: los LLM pueden descubrir y construir nuevos conceptos y nuevas herramientas—esta es una dirección totalmente nueva pero con gran potencial
\end{enumerate}
\section{AI for Math: razonamiento formal y LLM Agent}

\subsection{Limitaciones de los métodos puramente neuronales}

En años recientes los LLM han tenido un desempeño sobresaliente en competencias de matemáticas (por ejemplo, la mejora notable de O1 en el AMC 2024, el desempeño de nivel medalla de plata de AlphaProof en la IMO), pero los métodos puramente neuronales tienen debilidades fundamentales:

\begin{warningbox}{Dos problemas fundamentales de los modelos de lenguaje puros para hacer matemáticas}
\begin{enumerate}
    \item \textbf{Escasez de datos}: el razonamiento matemático riguroso requiere una gran cantidad de trayectorias de razonamiento de alta calidad o funciones de recompensa de alta calidad; más allá del nivel de competencias o de preparatoria, es difícil obtenerlas solo a partir de datos humanos
    \item \textbf{El razonamiento en lenguaje natural es difícil de verificar}: la corrección del razonamiento en lenguaje no puede comprobarse formalmente, lo cual es especialmente riesgoso en sistemas reales donde los casos límite son críticos
\end{enumerate}
\end{warningbox}

\subsection{Representación formal: una alternativa}

\begin{importantbox}{Idea central de los métodos formales}
\begin{enumerate}
    \item \textbf{Autoformalización} (Autoformalization): convertir enunciados matemáticos informales al lenguaje formal de sistemas como Lean, Coq o Isabelle
    \item \textbf{Demostrador neuronal} (Neural Prover): busca secuencias de estrategias de demostración (tactics) dentro de un entorno formal
    \item \textbf{Verificación}: una demostración formal puede comprobarse automáticamente mediante un asistente de demostración; esto resuelve el problema de la verificación
    \item \textbf{Datos sintéticos}: se pueden generar en gran cantidad demostraciones candidatas y obtener de inmediato retroalimentación sobre su corrección, lo que convierte el problema de demostración en una tarea similar a un juego
\end{enumerate}
\end{importantbox}

\subsubsection{Cómo funciona el asistente de demostración Lean}

Lean es una máquina de estados:
\begin{itemize}
    \item \textbf{Estado} (State): los objetivos (goals) que se necesitan demostrar en el momento actual y las hipótesis del contexto
    \item \textbf{Acción} (Tactic): reglas de simplificación que, al aplicarse, cambian el estado de la demostración
    \item \textbf{Objetivo}: encontrar una secuencia de tactics que resuelva todos los goals y llegue al estado QED
\end{itemize}

Por ejemplo, para demostrar «si $x$ es par, entonces $x^2$ es par»: el estado inicial contiene el objetivo $x^2 \bmod 2 = 0$; mediante tactics como \texttt{intro} y \texttt{rw} se simplifica paso a paso hasta llegar a $0 = 0$, es decir, QED.

\subsection{Copra: demostración de teoremas basada en LLM Agent}

\begin{importantbox}{Arquitectura del sistema Copra}
Copra es un Agent de demostración de teoremas basado en aprendizaje en contexto (In-Context Learning):
\begin{enumerate}
    \item Usa un LLM de vanguardia (como GPT-4) para predecir la siguiente tactic
    \item Proporciona el estado de la demostración y el historial de búsqueda al LLM como prompt
    \item Ejecuta la tactic predicha y obtiene el nuevo estado o retroalimentación de error
    \item Incorpora la información del error de vuelta al prompt para que el LLM corrija su estrategia
    \item En el nivel externo usa búsqueda en profundidad (DFS) para organizar todo el proceso de demostración
\end{enumerate}
\end{importantbox}

Ventajas centrales de Copra:
\begin{itemize}
    \item \textbf{No requiere entrenamiento adicional}: aprovecha directamente las capacidades de un LLM preentrenado
    \item \textbf{Se beneficia de inmediato de los avances de los LLM}: la actualización de GPT-4 a O1/O3 mejora el desempeño directamente
    \item \textbf{Combina el razonamiento en lenguaje natural con el razonamiento en lenguaje formal}
    \item \textbf{No requiere un corpus de entrenamiento}: es aplicable a problemas de investigación completamente nuevos
    \item \textbf{La retroalimentación del asistente de demostración provee un anclaje} (grounding): evita las alucinaciones
\end{itemize}

\subsection{Demostración jerárquica: el poder de la abstracción}

\begin{importantbox}{Descomposición jerárquica de demostraciones impulsada por LLM}
Para teoremas complejos (como los problemas de la IMO), Copra puede extenderse a una resolución jerárquica:
\begin{enumerate}
    \item Se le pide al LLM que genere un plan de demostración informal (informal proof plan)
    \item A partir del plan, se descompone el teorema en la formulación formal de varios subobjetivos (sub-goals)
    \item Se usa Copra para demostrar los subobjetivos uno por uno
    \item Se colocan los subobjetivos ya demostrados en el contexto y se demuestra el teorema completo
\end{enumerate}
Idea clave: el LLM no solo puede predecir pasos de demostración de bajo nivel, sino también realizar \textbf{razonamiento abstracto de alto nivel}; todo esto se logra mediante ingeniería de prompt, sin necesidad de entrenamiento adicional.
\end{importantbox}

Tomando como ejemplo el problema de la IMO «demostrar que $\frac{21n+4}{14n+3}$ es irreducible para todo número natural $n$»: O3 descompone el problema en tres lemas auxiliares sobre el GCD; después de que Copra los demuestra por separado, se combinan para demostrar el teorema completo.

\subsection{Aplicación: verificación formal de compiladores}

\begin{knowledgebox}{El valor práctico de la verificación formal}
La verificación formal puede demostrar matemáticamente la corrección, la seguridad y el desempeño de un sistema. El Departamento de Defensa de Estados Unidos llegó a desplegar en drones un núcleo Linux verificado formalmente (seL4), y un equipo rojo no pudo vulnerarlo en seis semanas.

Históricamente la verificación formal ha sido extremadamente costosa, pero el auge de AI for Math está \textbf{cambiando de manera fundamental la ecuación de costo-beneficio}.
\end{knowledgebox}

La clase mostró un caso sencillo de verificación de un compilador de expresiones aritméticas (en el lenguaje Coq):
\begin{itemize}
    \item Lenguaje fuente: expresiones aritméticas definidas de manera recursiva (suma, multiplicación)
    \item Lenguaje objetivo: una lista de instrucciones basada en una máquina de pila
    \item Teorema de corrección del compilador: el resultado de ejecutar el código compilado es igual al resultado de evaluar el lenguaje fuente
    \item Copra no puede demostrarlo de una sola vez, pero mediante el LLM se inventa un lema auxiliar (la corrección de la compilación de una sola instrucción) y, al combinarlos, se completa la demostración completa
\end{itemize}

\subsection{Resumen de la sección}

\begin{itemize}
    \item Los LLM Agent para el descubrimiento matemático son muy prometedores; los métodos de aprendizaje en contexto ya han mostrado capacidades sólidas
    \item La retroalimentación del asistente de demostración es fundamental para el anclaje (grounding)
    \item Un diseño de Agent avanzado puede combinar la formalización, la planificación jerárquica y la generación de demostraciones de bajo nivel
    \item Por ahora todavía se necesita una arquitectura de Agent; los métodos puramente de entrenamiento aún no alcanzan un nivel suficiente
\end{itemize}
\section{IA para la ciencia: descubrimiento científico guiado por LLM}

\subsection{De la regresión simbólica a la evolución guiada por LLM}

\begin{knowledgebox}{Problema de regresión simbólica}
La \textbf{regresión simbólica} (Symbolic Regression) es la tarea de descubrir ecuaciones a partir de datos observacionales empíricos. Por ejemplo, dados datos de medición del movimiento de Marte, descubrir la tercera ley de Kepler.

Los métodos tradicionales (como la programación genética PySR) mantienen una población de expresiones candidatas, que evoluciona mediante mutación y cruce. Los LLM pueden potenciar este proceso aportando conocimiento científico previo.
\end{knowledgebox}

\subsection{LaSR: regresión simbólica guiada por conceptos}

\begin{importantbox}{Marco LaSR (Language-guided Symbolic Regression)}
LaSR infiere conjuntamente hipótesis de programa y conceptos de alto nivel:

$$P(\pi, C \mid D) \propto P(D \mid \pi) \cdot P(\pi \mid C) \cdot P(C)$$

donde:
\begin{itemize}
    \item $P(D \mid \pi)$: la verosimilitud del programa $\pi$ respecto a los datos $D$ (calculada ejecutando el programa)
    \item $P(\pi \mid C)$: la probabilidad del programa dado el concepto $C$ (el conocimiento previo del LLM)
    \item $P(C)$: la distribución previa de los conceptos (el conocimiento científico del LLM)
\end{itemize}
\end{importantbox}

Flujo de trabajo de LaSR:
\begin{enumerate}
    \item Mantener una población de programas (ecuaciones) candidatos
    \item \textbf{Evolución de hipótesis}: mutación simbólica clásica + mutación guiada por LLM (basada en la biblioteca de conceptos)
    \item \textbf{Evaluación de aptitud}: qué tan bien se ajusta el programa a los datos
    \item \textbf{Abstracción de conceptos}: el LLM abstrae el conjunto de programas en descripciones de alto nivel en lenguaje natural (como «crecimiento/decaimiento exponencial», «ley de potencia»)
    \item \textbf{Evolución de conceptos}: el LLM combina conceptos existentes para generar conceptos nuevos (como «crecimiento exponencial» + «dependencia de la temperatura» $\to$ «distribución de Boltzmann»)
    \item La biblioteca de conceptos retroalimenta la evolución de hipótesis, formando un ciclo cerrado
\end{enumerate}

\subsection{Resultados experimentales}

\begin{itemize}
    \item En la tarea de descubrimiento de las ecuaciones de Feynman, LaSR supera tanto a la programación genética pura como a la versión sin aprendizaje de conceptos
    \item Incluso usando modelos de lenguaje locales (no de frontera), logra superar a la programación genética pura
    \item Los prompts del usuario (hints) pueden aportar una ganancia adicional
    \item La expresión de la ley de Coulomb que descubrió LaSR solo necesita 4 pasos de simplificación para llegar a la forma estándar, mientras que la expresión equivalente obtenida por PySR requiere una simplificación considerable
\end{itemize}

\subsection{Nuevo descubrimiento: leyes de escalamiento de LLM}

Para verificar que el método no se limita a «memorizar fórmulas conocidas», el equipo usó LaSR para descubrir una nueva ley de escalamiento de LLM que incorpora el parámetro del \textbf{número de ejemplos few-shot}:

\begin{warningbox}{Hallazgo interesante}
Un número grande de ejemplos few-shot resulta perjudicial para los modelos de baja capacidad, pero una vez que la capacidad del modelo supera cierto umbral, más shots producen un mejor desempeño. Al combinar este hallazgo con la ley de Chinchilla se obtuvo una ley de escalamiento más precisa.
\end{warningbox}

\subsection{Resumen de la sección}

\begin{itemize}
    \item La evolución guiada por LLM es una herramienta poderosa para la ciencia empírica
    \item La capacidad de abstracción del LLM (descubrimiento de conceptos) puede potenciar aún más el proceso de descubrimiento
    \item Se puede extender a entradas visuales de alta dimensión (usando VLM)
    \item Retos abiertos: validación de hipótesis, representación de conceptos más allá del lenguaje natural, y escalamiento a espacios de búsqueda mayores
\end{itemize}
\section{Síntesis y ampliación}

\subsection{Mensaje central}

Los Agentes de IA basados en LLM tienen un enorme potencial en el descubrimiento matemático y científico; su ventaja no radica solo en una capacidad de búsqueda sólida y en el conocimiento previo, sino sobre todo en la \textbf{capacidad de abstracción}: descubrir nuevos conceptos y construir soluciones jerárquicas. Los patrones de diseño clave incluyen:

\begin{enumerate}
    \item \textbf{LLM + verificación formal}: el LLM genera pasos de demostración y el asistente de demostración ofrece retroalimentación fundamentada
    \item \textbf{Descomposición jerárquica}: el LLM realiza la planeación de alto nivel y la descomposición en subproblemas
    \item \textbf{Evolución guiada por conceptos}: el LLM abstrae y combina conceptos de alto nivel para guiar la dirección de la búsqueda
    \item \textbf{Flexibilidad de la ingeniería de prompts}: permite lograr una mejora cualitativa de las capacidades sin entrenamiento adicional
\end{enumerate}

\subsection{Lecturas adicionales}
\begin{itemize}
    \item Copra: demostración de teoremas con Agentes de IA basados en LLM (Conference on Language Models, 2024)
    \item LaSR: regresión simbólica guiada por conceptos (NeurIPS 2024)
    \item FunSearch (DeepMind): búsqueda evolutiva guiada por LLM
    \item AlphaProof: RL + matemáticas formales
    \item Artículo de revisión ``AI for Math'' (con participación de Dawn Song, entre otros)
\end{itemize}

\end{document}
